% --- Archon physics formalization source begin ---
% archon:physics
% archon:covers IPhO2026Problems/problem_IPhO_2026_1_B_2.lean
% archon:source-report reports/ipho_2026/problem_IPhO_2026_1_B_2.source.json
% archon:problem-id IPhO_2026_1
% archon:part-id B.2

\chapter{Physics problem IPhO\_2026\_1\_B\_2}
\label{ch:IPhO2026Problems_problem_IPhO_2026_1_B_2}

\paragraph{Problem source.}
At one instant a positron and an electron, each of mass m and charges of equal
magnitude and opposite sign, are separated by 100*a\_0.  Their velocities are
antiparallel and perpendicular to their separation.  Each particle has angular
momentum of magnitude mu*hbar about the center of mass.  The system is isolated,
classical, non-relativistic, and has only electrostatic interaction.  The Bohr
radius is a\_0 = 4*pi*epsilon\_0*hbar\textasciicircum{}2/(m*e\textasciicircum{}2), and k = 1/(4*pi*epsilon\_0).

Current subquestion:
For mu = 15/2 the pair is unbound. Find the angle between the asymptotic relative velocity u\_infinity and the initial positron line of motion.

\paragraph{Current subquestion.}
For mu = 15/2 the pair is unbound. Find the angle between the asymptotic relative velocity u\_infinity and the initial positron line of motion.

\paragraph{Recorded answer/context.}
The signed deflection is -16.60 degrees, i.e. 16.60 degrees below the initial line of motion.

\paragraph{Figure/image path.}
/root/proposal\_for\_physic/science-mango/ipho\_2026\_source/image/T1\_page-2.png

\paragraph{Formalization target.}
create a compiling Lean file with sorry bodies at `IPhO2026Problems/problem\_IPhO\_2026\_1\_B\_2.lean`. The Lean declarations must preserve the physical quantities, dimensions or dimensional roles, figure labels, governing-law hypotheses, and final relation expressed by this problem.
Use Mathlib/Physlib names found through LeanExplore where available. If a physics API is missing, introduce faithful local abstractions rather than scalar placeholder aliases.

\begin{definition}[Plane]
  \label{decl:physics:IPhO_2026_1_B_2:Plane}
  \lean{IPhO2026Problems.IPhO2026_1_B_2.Plane}
  The oriented two-dimensional plane containing the trajectories in Figure 1b.
\end{definition}
\begin{proof}
  This is the defining declaration; unfolding it gives the stated typed data or relation.
\end{proof}

\begin{definition}[velocityDimension]
  \label{decl:physics:IPhO_2026_1_B_2:velocityDimension}
  \lean{IPhO2026Problems.IPhO2026_1_B_2.velocityDimension}
  Velocity dimension “L T⁻¹”.
\end{definition}
\begin{proof}
  This is the defining declaration; unfolding it gives the stated typed data or relation.
\end{proof}

\begin{definition}[angularMomentumDimension]
  \label{decl:physics:IPhO_2026_1_B_2:angularMomentumDimension}
  \lean{IPhO2026Problems.IPhO2026_1_B_2.angularMomentumDimension}
  Angular-momentum/action dimension “M L² T⁻¹”.
\end{definition}
\begin{proof}
  This is the defining declaration; unfolding it gives the stated typed data or relation.
\end{proof}

\begin{definition}[permittivityDimension]
  \label{decl:physics:IPhO_2026_1_B_2:permittivityDimension}
  \lean{IPhO2026Problems.IPhO2026_1_B_2.permittivityDimension}
  Vacuum-permittivity dimension “C² T² M⁻¹ L⁻³”.
\end{definition}
\begin{proof}
  This is the defining declaration; unfolding it gives the stated typed data or relation.
\end{proof}

\begin{definition}[coulombConstantDimension]
  \label{decl:physics:IPhO_2026_1_B_2:coulombConstantDimension}
  \lean{IPhO2026Problems.IPhO2026_1_B_2.coulombConstantDimension}
  Coulomb-constant dimension “M L³ T⁻² C⁻²”.
\end{definition}
\begin{proof}
  This is the defining declaration; unfolding it gives the stated typed data or relation.
\end{proof}

\begin{definition}[DimLength]
  \label{decl:physics:IPhO_2026_1_B_2:DimLength}
  \lean{IPhO2026Problems.IPhO2026_1_B_2.DimLength}
  The local abbreviation DimLength records part of the typed physics model.
\end{definition}
\begin{proof}
  This is the defining declaration; unfolding it gives the stated typed data or relation.
\end{proof}

\begin{definition}[DimMass]
  \label{decl:physics:IPhO_2026_1_B_2:DimMass}
  \lean{IPhO2026Problems.IPhO2026_1_B_2.DimMass}
  The local abbreviation DimMass records part of the typed physics model.
\end{definition}
\begin{proof}
  This is the defining declaration; unfolding it gives the stated typed data or relation.
\end{proof}

\begin{definition}[DimCharge]
  \label{decl:physics:IPhO_2026_1_B_2:DimCharge}
  \lean{IPhO2026Problems.IPhO2026_1_B_2.DimCharge}
  The local abbreviation DimCharge records part of the typed physics model.
\end{definition}
\begin{proof}
  This is the defining declaration; unfolding it gives the stated typed data or relation.
\end{proof}

\begin{definition}[DimPosition]
  \label{decl:physics:IPhO_2026_1_B_2:DimPosition}
  \lean{IPhO2026Problems.IPhO2026_1_B_2.DimPosition}
  \uses{decl:physics:IPhO_2026_1_B_2:Plane}
  The local abbreviation DimPosition records part of the typed physics model.
\end{definition}
\begin{proof}
  This is the defining declaration; unfolding it gives the stated typed data or relation.
\end{proof}

\begin{definition}[DimVelocityVector]
  \label{decl:physics:IPhO_2026_1_B_2:DimVelocityVector}
  \lean{IPhO2026Problems.IPhO2026_1_B_2.DimVelocityVector}
  \uses{decl:physics:IPhO_2026_1_B_2:Plane, decl:physics:IPhO_2026_1_B_2:velocityDimension}
  The local abbreviation DimVelocityVector records part of the typed physics model.
\end{definition}
\begin{proof}
  This is the defining declaration; unfolding it gives the stated typed data or relation.
\end{proof}

\begin{definition}[DimSpeed]
  \label{decl:physics:IPhO_2026_1_B_2:DimSpeed}
  \lean{IPhO2026Problems.IPhO2026_1_B_2.DimSpeed}
  \uses{decl:physics:IPhO_2026_1_B_2:velocityDimension}
  The local abbreviation DimSpeed records part of the typed physics model.
\end{definition}
\begin{proof}
  This is the defining declaration; unfolding it gives the stated typed data or relation.
\end{proof}

\begin{definition}[DimAngularMomentum]
  \label{decl:physics:IPhO_2026_1_B_2:DimAngularMomentum}
  \lean{IPhO2026Problems.IPhO2026_1_B_2.DimAngularMomentum}
  \uses{decl:physics:IPhO_2026_1_B_2:angularMomentumDimension}
  The local abbreviation DimAngularMomentum records part of the typed physics model.
\end{definition}
\begin{proof}
  This is the defining declaration; unfolding it gives the stated typed data or relation.
\end{proof}

\begin{definition}[DimPermittivity]
  \label{decl:physics:IPhO_2026_1_B_2:DimPermittivity}
  \lean{IPhO2026Problems.IPhO2026_1_B_2.DimPermittivity}
  \uses{decl:physics:IPhO_2026_1_B_2:permittivityDimension}
  The local abbreviation DimPermittivity records part of the typed physics model.
\end{definition}
\begin{proof}
  This is the defining declaration; unfolding it gives the stated typed data or relation.
\end{proof}

\begin{definition}[DimCoulombConstant]
  \label{decl:physics:IPhO_2026_1_B_2:DimCoulombConstant}
  \lean{IPhO2026Problems.IPhO2026_1_B_2.DimCoulombConstant}
  \uses{decl:physics:IPhO_2026_1_B_2:coulombConstantDimension}
  The local abbreviation DimCoulombConstant records part of the typed physics model.
\end{definition}
\begin{proof}
  This is the defining declaration; unfolding it gives the stated typed data or relation.
\end{proof}

\begin{definition}[scalarSI]
  \label{decl:physics:IPhO_2026_1_B_2:scalarSI}
  \lean{IPhO2026Problems.IPhO2026_1_B_2.scalarSI}
  The scalar SI readout of a dimensionful scalar quantity.
\end{definition}
\begin{proof}
  This is the defining declaration; unfolding it gives the stated typed data or relation.
\end{proof}

\begin{definition}[vectorSI]
  \label{decl:physics:IPhO_2026_1_B_2:vectorSI}
  \lean{IPhO2026Problems.IPhO2026_1_B_2.vectorSI}
  \uses{decl:physics:IPhO_2026_1_B_2:Plane}
  The component vector in the standard SI units for its dimension.
\end{definition}
\begin{proof}
  This is the defining declaration; unfolding it gives the stated typed data or relation.
\end{proof}

\begin{definition}[ParticleLabel]
  \label{decl:physics:IPhO_2026_1_B_2:ParticleLabel}
  \lean{IPhO2026Problems.IPhO2026_1_B_2.ParticleLabel}
  \uses{decl:physics:IPhO_2026_1_B_2:Figure1bInitialConditions}
  The two labeled particles in Figure 1b.
\end{definition}
\begin{proof}
  This is the defining declaration; unfolding it gives the stated typed data or relation.
\end{proof}

\begin{definition}[PhysicalConstants]
  \label{decl:physics:IPhO_2026_1_B_2:PhysicalConstants}
  \lean{IPhO2026Problems.IPhO2026_1_B_2.PhysicalConstants}
  \uses{decl:physics:IPhO_2026_1_B_2:DimLength, decl:physics:IPhO_2026_1_B_2:DimMass, decl:physics:IPhO_2026_1_B_2:DimCharge, decl:physics:IPhO_2026_1_B_2:DimSpeed, decl:physics:IPhO_2026_1_B_2:DimAngularMomentum, decl:physics:IPhO_2026_1_B_2:DimPermittivity, decl:physics:IPhO_2026_1_B_2:DimCoulombConstant, decl:physics:IPhO_2026_1_B_2:scalarSI}
  The universal constants and common particle mass used in the problem.
\end{definition}
\begin{proof}
  This is the defining declaration; unfolding it gives the stated typed data or relation.
\end{proof}

\begin{definition}[ConstantRelations]
  \label{decl:physics:IPhO_2026_1_B_2:ConstantRelations}
  \lean{IPhO2026Problems.IPhO2026_1_B_2.ConstantRelations}
  \uses{decl:physics:IPhO_2026_1_B_2:scalarSI, decl:physics:IPhO_2026_1_B_2:PhysicalConstants}
  The two constant identities printed on the source page: “a₀ = 4 π ε₀ ℏ² / (m e²)” and “k = 1 / (4 π ε₀)”.
\end{definition}
\begin{proof}
  This is the defining declaration; unfolding it gives the stated typed data or relation.
\end{proof}

\begin{definition}[PairMotion]
  \label{decl:physics:IPhO_2026_1_B_2:PairMotion}
  \lean{IPhO2026Problems.IPhO2026_1_B_2.PairMotion}
  \uses{decl:physics:IPhO_2026_1_B_2:DimMass, decl:physics:IPhO_2026_1_B_2:DimCharge, decl:physics:IPhO_2026_1_B_2:DimPosition, decl:physics:IPhO_2026_1_B_2:DimVelocityVector, decl:physics:IPhO_2026_1_B_2:ParticleLabel}
  The dimensionful positions and velocities of both particles.
\end{definition}
\begin{proof}
  This is the defining declaration; unfolding it gives the stated typed data or relation.
\end{proof}

\begin{definition}[positionSI]
  \label{decl:physics:IPhO_2026_1_B_2:positionSI}
  \lean{IPhO2026Problems.IPhO2026_1_B_2.positionSI}
  \uses{decl:physics:IPhO_2026_1_B_2:Plane, decl:physics:IPhO_2026_1_B_2:vectorSI, decl:physics:IPhO_2026_1_B_2:ParticleLabel, decl:physics:IPhO_2026_1_B_2:PairMotion}
  SI position vector of a labeled particle at the specified time in seconds.
\end{definition}
\begin{proof}
  This is the defining declaration; unfolding it gives the stated typed data or relation.
\end{proof}

\begin{definition}[velocitySI]
  \label{decl:physics:IPhO_2026_1_B_2:velocitySI}
  \lean{IPhO2026Problems.IPhO2026_1_B_2.velocitySI}
  \uses{decl:physics:IPhO_2026_1_B_2:Plane, decl:physics:IPhO_2026_1_B_2:vectorSI, decl:physics:IPhO_2026_1_B_2:ParticleLabel, decl:physics:IPhO_2026_1_B_2:PairMotion}
  SI velocity vector of a labeled particle at the specified time in seconds.
\end{definition}
\begin{proof}
  This is the defining declaration; unfolding it gives the stated typed data or relation.
\end{proof}

\begin{definition}[particleMassSI]
  \label{decl:physics:IPhO_2026_1_B_2:particleMassSI}
  \lean{IPhO2026Problems.IPhO2026_1_B_2.particleMassSI}
  \uses{decl:physics:IPhO_2026_1_B_2:scalarSI, decl:physics:IPhO_2026_1_B_2:ParticleLabel, decl:physics:IPhO_2026_1_B_2:PairMotion}
  SI mass readout of a labeled particle.
\end{definition}
\begin{proof}
  This is the defining declaration; unfolding it gives the stated typed data or relation.
\end{proof}

\begin{definition}[particleChargeSI]
  \label{decl:physics:IPhO_2026_1_B_2:particleChargeSI}
  \lean{IPhO2026Problems.IPhO2026_1_B_2.particleChargeSI}
  \uses{decl:physics:IPhO_2026_1_B_2:scalarSI, decl:physics:IPhO_2026_1_B_2:ParticleLabel, decl:physics:IPhO_2026_1_B_2:PairMotion}
  SI charge readout of a labeled particle.
\end{definition}
\begin{proof}
  This is the defining declaration; unfolding it gives the stated typed data or relation.
\end{proof}

\begin{definition}[relativeDisplacementSI]
  \label{decl:physics:IPhO_2026_1_B_2:relativeDisplacementSI}
  \lean{IPhO2026Problems.IPhO2026_1_B_2.relativeDisplacementSI}
  \uses{decl:physics:IPhO_2026_1_B_2:Plane, decl:physics:IPhO_2026_1_B_2:PairMotion, decl:physics:IPhO_2026_1_B_2:positionSI}
  Relative displacement from the electron to the positron, in metres.
\end{definition}
\begin{proof}
  This is the defining declaration; unfolding it gives the stated typed data or relation.
\end{proof}

\begin{definition}[relativeVelocitySI]
  \label{decl:physics:IPhO_2026_1_B_2:relativeVelocitySI}
  \lean{IPhO2026Problems.IPhO2026_1_B_2.relativeVelocitySI}
  \uses{decl:physics:IPhO_2026_1_B_2:Plane, decl:physics:IPhO_2026_1_B_2:PairMotion, decl:physics:IPhO_2026_1_B_2:velocitySI}
  Positron velocity relative to the electron, in metres per second.
\end{definition}
\begin{proof}
  This is the defining declaration; unfolding it gives the stated typed data or relation.
\end{proof}

\begin{definition}[separationSI]
  \label{decl:physics:IPhO_2026_1_B_2:separationSI}
  \lean{IPhO2026Problems.IPhO2026_1_B_2.separationSI}
  \uses{decl:physics:IPhO_2026_1_B_2:PairMotion, decl:physics:IPhO_2026_1_B_2:relativeDisplacementSI}
  Electron--positron separation in metres.
\end{definition}
\begin{proof}
  This is the defining declaration; unfolding it gives the stated typed data or relation.
\end{proof}

\begin{definition}[Figure1bFrame]
  \label{decl:physics:IPhO_2026_1_B_2:Figure1bFrame}
  \lean{IPhO2026Problems.IPhO2026_1_B_2.Figure1bFrame}
  \uses{decl:physics:IPhO_2026_1_B_2:Plane}
  The oriented axes read from Figure 1b.
\end{definition}
\begin{proof}
  This is the defining declaration; unfolding it gives the stated typed data or relation.
\end{proof}

\begin{definition}[signedAngularMomentumSI]
  \label{decl:physics:IPhO_2026_1_B_2:signedAngularMomentumSI}
  \lean{IPhO2026Problems.IPhO2026_1_B_2.signedAngularMomentumSI}
  \uses{decl:physics:IPhO_2026_1_B_2:DimPosition, decl:physics:IPhO_2026_1_B_2:vectorSI, decl:physics:IPhO_2026_1_B_2:ParticleLabel, decl:physics:IPhO_2026_1_B_2:PairMotion, decl:physics:IPhO_2026_1_B_2:positionSI, decl:physics:IPhO_2026_1_B_2:velocitySI, decl:physics:IPhO_2026_1_B_2:particleMassSI, decl:physics:IPhO_2026_1_B_2:Figure1bFrame}
  Signed angular momentum of one particle about a chosen center, in SI units.
\end{definition}
\begin{proof}
  This is the defining declaration; unfolding it gives the stated typed data or relation.
\end{proof}

\begin{definition}[totalAngularMomentumMagnitudeSI]
  \label{decl:physics:IPhO_2026_1_B_2:totalAngularMomentumMagnitudeSI}
  \lean{IPhO2026Problems.IPhO2026_1_B_2.totalAngularMomentumMagnitudeSI}
  \uses{decl:physics:IPhO_2026_1_B_2:DimPosition, decl:physics:IPhO_2026_1_B_2:PairMotion, decl:physics:IPhO_2026_1_B_2:Figure1bFrame, decl:physics:IPhO_2026_1_B_2:signedAngularMomentumSI}
  Magnitude of the pair's total angular momentum about the center of mass.
\end{definition}
\begin{proof}
  This is the defining declaration; unfolding it gives the stated typed data or relation.
\end{proof}

\begin{definition}[Figure1bInitialConditions]
  \label{decl:physics:IPhO_2026_1_B_2:Figure1bInitialConditions}
  \lean{IPhO2026Problems.IPhO2026_1_B_2.Figure1bInitialConditions}
  \uses{decl:physics:IPhO_2026_1_B_2:DimPosition, decl:physics:IPhO_2026_1_B_2:DimSpeed, decl:physics:IPhO_2026_1_B_2:scalarSI, decl:physics:IPhO_2026_1_B_2:vectorSI, decl:physics:IPhO_2026_1_B_2:PhysicalConstants, decl:physics:IPhO_2026_1_B_2:PairMotion, decl:physics:IPhO_2026_1_B_2:positionSI, decl:physics:IPhO_2026_1_B_2:velocitySI, decl:physics:IPhO_2026_1_B_2:particleMassSI, decl:physics:IPhO_2026_1_B_2:particleChargeSI, decl:physics:IPhO_2026_1_B_2:relativeDisplacementSI, decl:physics:IPhO_2026_1_B_2:Figure1bFrame, decl:physics:IPhO_2026_1_B_2:signedAngularMomentumSI}
  All numerical and geometric data at the pictured instant.
\end{definition}
\begin{proof}
  This is the defining declaration; unfolding it gives the stated typed data or relation.
\end{proof}

\begin{definition}[coulombForceOnPositronSI]
  \label{decl:physics:IPhO_2026_1_B_2:coulombForceOnPositronSI}
  \lean{IPhO2026Problems.IPhO2026_1_B_2.coulombForceOnPositronSI}
  \uses{decl:physics:IPhO_2026_1_B_2:Plane, decl:physics:IPhO_2026_1_B_2:scalarSI, decl:physics:IPhO_2026_1_B_2:PhysicalConstants, decl:physics:IPhO_2026_1_B_2:PairMotion, decl:physics:IPhO_2026_1_B_2:particleChargeSI, decl:physics:IPhO_2026_1_B_2:relativeDisplacementSI}
  Coulomb force on the positron, evaluated in SI components.
\end{definition}
\begin{proof}
  This is the defining declaration; unfolding it gives the stated typed data or relation.
\end{proof}

\begin{definition}[CoulombDynamics]
  \label{decl:physics:IPhO_2026_1_B_2:CoulombDynamics}
  \lean{IPhO2026Problems.IPhO2026_1_B_2.CoulombDynamics}
  \uses{decl:physics:IPhO_2026_1_B_2:Plane, decl:physics:IPhO_2026_1_B_2:scalarSI, decl:physics:IPhO_2026_1_B_2:ParticleLabel, decl:physics:IPhO_2026_1_B_2:PhysicalConstants, decl:physics:IPhO_2026_1_B_2:PairMotion, decl:physics:IPhO_2026_1_B_2:positionSI, decl:physics:IPhO_2026_1_B_2:velocitySI, decl:physics:IPhO_2026_1_B_2:particleMassSI, decl:physics:IPhO_2026_1_B_2:relativeDisplacementSI, decl:physics:IPhO_2026_1_B_2:coulombForceOnPositronSI}
  The classical non-relativistic, isolated two-body model.
\end{definition}
\begin{proof}
  This is the defining declaration; unfolding it gives the stated typed data or relation.
\end{proof}

\begin{definition}[initialTotalEnergySI]
  \label{decl:physics:IPhO_2026_1_B_2:initialTotalEnergySI}
  \lean{IPhO2026Problems.IPhO2026_1_B_2.initialTotalEnergySI}
  \uses{decl:physics:IPhO_2026_1_B_2:scalarSI, decl:physics:IPhO_2026_1_B_2:PhysicalConstants, decl:physics:IPhO_2026_1_B_2:PairMotion, decl:physics:IPhO_2026_1_B_2:velocitySI, decl:physics:IPhO_2026_1_B_2:particleMassSI, decl:physics:IPhO_2026_1_B_2:particleChargeSI, decl:physics:IPhO_2026_1_B_2:separationSI}
  Total mechanical energy computed at the pictured instant, in joules.
\end{definition}
\begin{proof}
  This is the defining declaration; unfolding it gives the stated typed data or relation.
\end{proof}

\begin{definition}[ConicOrbitData]
  \label{decl:physics:IPhO_2026_1_B_2:ConicOrbitData}
  \lean{IPhO2026Problems.IPhO2026_1_B_2.ConicOrbitData}
  \uses{decl:physics:IPhO_2026_1_B_2:Plane, decl:physics:IPhO_2026_1_B_2:DimLength, decl:physics:IPhO_2026_1_B_2:DimAngularMomentum}
  Data used by the two conic-orbit hints on the official source page.
\end{definition}
\begin{proof}
  This is the defining declaration; unfolding it gives the stated typed data or relation.
\end{proof}

\begin{definition}[ConicOrbitLaws]
  \label{decl:physics:IPhO_2026_1_B_2:ConicOrbitLaws}
  \lean{IPhO2026Problems.IPhO2026_1_B_2.ConicOrbitLaws}
  \uses{decl:physics:IPhO_2026_1_B_2:scalarSI, decl:physics:IPhO_2026_1_B_2:PhysicalConstants, decl:physics:IPhO_2026_1_B_2:PairMotion, decl:physics:IPhO_2026_1_B_2:relativeDisplacementSI, decl:physics:IPhO_2026_1_B_2:separationSI, decl:physics:IPhO_2026_1_B_2:Figure1bFrame, decl:physics:IPhO_2026_1_B_2:totalAngularMomentumMagnitudeSI, decl:physics:IPhO_2026_1_B_2:Figure1bInitialConditions, decl:physics:IPhO_2026_1_B_2:initialTotalEnergySI, decl:physics:IPhO_2026_1_B_2:ConicOrbitData}
  The supplied Kepler/Coulomb eccentricity law and polar conic equation.
\end{definition}
\begin{proof}
  This is the defining declaration; unfolding it gives the stated typed data or relation.
\end{proof}

\begin{definition}[IsUnbound]
  \label{decl:physics:IPhO_2026_1_B_2:IsUnbound}
  \lean{IPhO2026Problems.IPhO2026_1_B_2.IsUnbound}
  \uses{decl:physics:IPhO_2026_1_B_2:PairMotion, decl:physics:IPhO_2026_1_B_2:separationSI}
  The stated unbound condition: the separation tends to infinity in the future.
\end{definition}
\begin{proof}
  This is the defining declaration; unfolding it gives the stated typed data or relation.
\end{proof}

\begin{definition}[signedDeflectionDegrees]
  \label{decl:physics:IPhO_2026_1_B_2:signedDeflectionDegrees}
  \lean{IPhO2026Problems.IPhO2026_1_B_2.signedDeflectionDegrees}
  \uses{decl:physics:IPhO_2026_1_B_2:Plane, decl:physics:IPhO_2026_1_B_2:PairMotion, decl:physics:IPhO_2026_1_B_2:velocitySI, decl:physics:IPhO_2026_1_B_2:Figure1bFrame}
  The signed counterclockwise angle from the initial positron velocity to the asymptotic positron-relative-to-electron velocity, reported in degrees.
\end{definition}
\begin{proof}
  This is the defining declaration; unfolding it gives the stated typed data or relation.
\end{proof}

\begin{definition}[RoundsToNearestHundredth]
  \label{decl:physics:IPhO_2026_1_B_2:RoundsToNearestHundredth}
  \lean{IPhO2026Problems.IPhO2026_1_B_2.RoundsToNearestHundredth}
  “actual” rounds to “reported” at two digits after the decimal point.
\end{definition}
\begin{proof}
  This is the defining declaration; unfolding it gives the stated typed data or relation.
\end{proof}

\begin{theorem}[Physics formalization target]
\label{thm:physics:IPhO_2026_1_B_2:target}
\lean{IPhO2026Problems.IPhO2026_1_B_2.IPhO_2026_1_B_2}
\uses{decl:physics:IPhO_2026_1_B_2:Plane, decl:physics:IPhO_2026_1_B_2:velocityDimension, decl:physics:IPhO_2026_1_B_2:angularMomentumDimension, decl:physics:IPhO_2026_1_B_2:permittivityDimension, decl:physics:IPhO_2026_1_B_2:coulombConstantDimension, decl:physics:IPhO_2026_1_B_2:DimLength, decl:physics:IPhO_2026_1_B_2:DimMass, decl:physics:IPhO_2026_1_B_2:DimCharge, decl:physics:IPhO_2026_1_B_2:DimPosition, decl:physics:IPhO_2026_1_B_2:DimVelocityVector, decl:physics:IPhO_2026_1_B_2:DimSpeed, decl:physics:IPhO_2026_1_B_2:DimAngularMomentum, decl:physics:IPhO_2026_1_B_2:DimPermittivity, decl:physics:IPhO_2026_1_B_2:DimCoulombConstant, decl:physics:IPhO_2026_1_B_2:scalarSI, decl:physics:IPhO_2026_1_B_2:vectorSI, decl:physics:IPhO_2026_1_B_2:ParticleLabel, decl:physics:IPhO_2026_1_B_2:PhysicalConstants, decl:physics:IPhO_2026_1_B_2:ConstantRelations, decl:physics:IPhO_2026_1_B_2:PairMotion, decl:physics:IPhO_2026_1_B_2:positionSI, decl:physics:IPhO_2026_1_B_2:velocitySI, decl:physics:IPhO_2026_1_B_2:particleMassSI, decl:physics:IPhO_2026_1_B_2:particleChargeSI, decl:physics:IPhO_2026_1_B_2:relativeDisplacementSI, decl:physics:IPhO_2026_1_B_2:relativeVelocitySI, decl:physics:IPhO_2026_1_B_2:separationSI, decl:physics:IPhO_2026_1_B_2:Figure1bFrame, decl:physics:IPhO_2026_1_B_2:signedAngularMomentumSI, decl:physics:IPhO_2026_1_B_2:totalAngularMomentumMagnitudeSI, decl:physics:IPhO_2026_1_B_2:Figure1bInitialConditions, decl:physics:IPhO_2026_1_B_2:coulombForceOnPositronSI, decl:physics:IPhO_2026_1_B_2:CoulombDynamics, decl:physics:IPhO_2026_1_B_2:initialTotalEnergySI, decl:physics:IPhO_2026_1_B_2:ConicOrbitData, decl:physics:IPhO_2026_1_B_2:ConicOrbitLaws, decl:physics:IPhO_2026_1_B_2:IsUnbound, decl:physics:IPhO_2026_1_B_2:signedDeflectionDegrees, decl:physics:IPhO_2026_1_B_2:RoundsToNearestHundredth}
The assigned autoformalize agent should translate this physics problem into Lean declarations in the covered file, with theorem and lemma proof bodies written as `by sorry`.
\end{theorem}
\begin{proof}
This is an autoformalization task, not a proof task. Produce statements that can later be proved by the physics prover without weakening the physical model.
\end{proof}
% --- Archon physics formalization source end ---
